%% sections/03-pipeline.tex

\section{The Amendment Pipeline}
\label{sec:pipeline}

\subsection{Stage 1: Innovation Logging}

An innovation log entry is created immediately after the session-gate scoring stage
when the gate reviewer identifies an outcome that:

\begin{itemize}
  \item is clearly a quality achievement (not a failure);
  \item is not described by any existing rubric anchor;
  \item could plausibly recur in future sessions (not purely session-specific).
\end{itemize}

Each entry has five fields:

\begin{description}
  \item[Technique:] A short title naming the pattern (e.g., ``Cross-character
    mechanical compound,'' ``Sustained behavioral pattern as choice'').
  \item[Dimension:] The dimension the innovation is most relevant to (one primary;
    cross-dimensions noted as secondary).
  \item[What happened:] A direct quotation from the session log with scene reference,
    plus a brief description of why this exceeded the rubric's prediction.
  \item[Scope:] Party-specific (likely to recur with this party only), adventure-specific
    (likely to recur with this adventure design pattern), or universal (likely to recur
    with any party or adventure).
  \item[Status:] logged $\to$ proposed $\to$ adopted (vX.Y) or reviewed-not-adopted.
\end{description}

The entries are append-only: no prior entry is deleted or modified after creation.
Status updates are the only modifications permitted.

\subsection{Stage 2: Clustering}

After each session, the full log of all entries with status ``logged'' is scanned for
dimension clusters. The threshold is:

\[
  \text{Trigger} = |\{e \in \text{log} : e.\text{dimension} = d \land e.\text{status} = \text{logged}\}| \geq 2
\]

for any dimension $d$. When the threshold fires, an amendment proposal is drafted.

Clustering does not require entries to be semantically identical; it requires that
they address the same dimension from different instances. In practice, clusters are
thematically coherent: entries that co-cluster address the same underlying quality
phenomenon, observed from different session contexts. Table~\ref{tab:cluster-examples}
shows the constituent innovations for three representative amendments.

\begin{table*}[t]
\caption{Representative Amendment Clusters: Constituent Innovations and Resulting Anchor}
\label{tab:cluster-examples}
\begin{tabular}{@{}llp{5cm}p{5cm}@{}}
\toprule
Amendment & Sessions & Constituent Innovations (abbreviated) & Resulting Anchor Text \\
\midrule
v1.1 AL & S01 & I-S01-02 (Grom receives atmos flag Kessa found); I-S01-10 (Grom sees fresco, Aelric fumbled); I-S01-12 (grief-chains without mechanics) & Per-PC two-column finder/receiver tracking; inter-PC chain-reaction at 10-band \\
v1.4 MF & S03–S04 & I-S03-03 (Portent replaces Religion, framed in-voice); I-S04-01 (same pattern, second instance) & Cross-character mechanical compound at 10-band; Portent-as-benediction as exemplar \\
v1.9 CA & C4-S1 & I-C4-S1-02 (Tomek token earned through non-credit-seeking pattern); I-C4-S1-04 (four PCs independently withhold Luca disclosure) & Sustained behavioral pattern at 9+; independent-motivation convergence as second path \\
\bottomrule
\end{tabular}
\end{table*}

\subsection{Stage 3: Amendment Adoption}

An amendment proposal is reviewed against two criteria: (a) the described pattern is
not already captured by an existing anchor, and (b) the proposed anchor text is
specific enough to produce consistent scores on future sessions where the pattern appears.

Proposals that pass review are adopted: the rubric version is incremented, constituent
innovations are marked adopted (vX.Y), and the new anchor is added to the rubric
document. Proposals that fail review are marked reviewed-not-adopted with a brief
reason. In the Marathon corpus, all proposals were adopted; no proposals were
rejected. This may reflect the conservative review standard --- proposals were only
drafted when the pattern was clearly novel and had strong textual evidence.

\subsection{Forward-Only Constraint}

The forward-only constraint means that sessions scored before an amendment cannot be
re-scored under the new anchor. This creates version-windows: sessions S01--S02
were scored under v1.0; sessions S03 under v1.1--v1.2; etc. Cross-version comparison
is valid when the applicable version is treated as an explicit covariate, and invalid
when it is ignored.

The forward-only constraint serves a practical function beyond statistical cleanliness:
it prevents retroactive inflation. If amendments could be applied retroactively, every
improvement in rubric expressiveness would immediately improve all prior scores,
obscuring whether actual session quality changed or just the measurement changed.
Forward-only application separates measurement improvement from quality improvement.
